\chapter{El oscilador armónico cuántico}\label{ch:b3:harmonic-oscillator}

Casi nada en la naturaleza es exactamente un oscilador armónico, y casi todo
lo es de forma aproximada: cualquier sistema apartado de un equilibrio
estable --- el enlace de una molécula, un átomo en un cristal, un puente, un
modo del campo electromagnético --- siente una fuerza recuperadora
proporcional al desplazamiento, porque todo potencial suave es una parábola
en el fondo de su pozo. Quien resuelve una vez el oscilador cuántico
resuelve, por tanto, las pequeñas vibraciones del mundo entero. Este
capítulo lo resuelve al estilo algebraico que se ha convertido en la firma
de la mecánica cuántica: dos \omterm{def:b3:harmonic-oscillator:ladder}{operadores escalera} suben y bajan una escalera
perfectamente regular de niveles $\big(n +
\tfrac12\big)\hbar\omega$, siendo el medio peldaño del fondo --- la \omterm{thm:b3:harmonic-oscillator:spectrum}{energía
del punto cero} --- un teorema y no una opción. Las consecuencias van desde
por qué el helio no se congela nunca hasta cómo una molécula de dióxido de
carbono, resonando en su propio $\hbar\omega$, intercepta el calor que la
Tierra emite al espacio.

\section{La escalera}

\begin{definition}[Operadores escalera]\label{def:b3:harmonic-oscillator:ladder}
\index{operadores escalera}\index{operador de aniquilación}\index{operador de creación}
Para $\hat H = \hat p^2/2m + \tfrac12 m\omega^2\hat x^2$, introdúzcase la
pareja adimensional y no hermítica
\[
\hat a = \sqrt{\frac{m\omega}{2\hbar}}\Big(\hat x +
\frac{\iu\hat p}{m\omega}\Big) , \qquad
\hat a^\dagger = \sqrt{\frac{m\omega}{2\hbar}}\Big(\hat x -
\frac{\iu\hat p}{m\omega}\Big) .
\]
A partir de $[\hat x, \hat p] = \iu\hbar$:
\[
[\hat a, \hat a^\dagger] = 1 , \qquad
\hat H = \hbar\omega\Big(\hat N + \tfrac12\Big) , \quad
\hat N = \hat a^\dagger\hat a .
\]
$\hat N$ es hermítico; sus \omterm{def:b3:quantum-formalism:observable}{valores propios} contarán \emph{cuantos}, y
$\hat a$ y $\hat a^\dagger$ --- los operadores de \emph{aniquilación} y de
\emph{creación} --- quitarán y añadirán uno.
\end{definition}

\begin{theorem}[El espectro, solo con álgebra]\label{thm:b3:harmonic-oscillator:spectrum}
\index{energía del punto cero}
Los \omterm{def:b3:quantum-formalism:observable}{valores propios} de $\hat H$ son exactamente
\[
E_n = \Big(n + \tfrac12\Big)\hbar\omega , \qquad n = 0, 1, 2, \dots
\]
--- una escalera infinita de peldaños iguales $\hbar\omega$ sobre un nivel
fundamental que \emph{no} es nulo: la \emph{\omterm{thm:b3:harmonic-oscillator:spectrum}{energía del punto cero}}
$\tfrac12\hbar\omega$. Los estados propios normalizados $\ket n$ están
conectados por
\[
\hat a\ket n = \sqrt n\,\ket{n - 1} , \qquad
\hat a^\dagger\ket n = \sqrt{n + 1}\,\ket{n + 1} ,
\qquad
\ket n = \frac{(\hat a^\dagger)^n}{\sqrt{n!}}\,\ket0 .
\]
\end{theorem}

\begin{proof}
A partir del \omterm{def:b3:quantum-formalism:commutator}{conmutador}, $\hat N\hat a = \hat a(\hat N - 1)$: si
$\ket\nu$ es vector propio de $\hat N$ con \omterm{def:b3:quantum-formalism:observable}{valor propio} $\nu$, entonces $\hat a\ket\nu$ es
un vector propio con $\nu - 1$ (o bien el vector nulo). Cada descenso solo
está permitido mientras $\nu \ge 0$, ya que $\nu = \braket{\nu}{\hat
N\nu} = \|\hat a\ket\nu\|^2 \ge 0$; el descenso debe, pues, terminar, y solo
termina en un estado con $\hat a\ket{\nu_0}
= 0$, de donde $\nu_0 = 0$. Así que $\nu$ recorre los enteros no
negativos. Las normalizaciones se siguen de $\|\hat a\ket n\|^2 = n$
y $\|\hat a^\dagger\ket n\|^2 = n + 1$.
\end{proof}

\begin{omfigure}
\begin{tikzpicture}[font=\small, >=Stealth]
  \draw[->] (0,-0.3) -- (0,5.4);
  \node[anchor=south] at (0,5.4) {$E$};
  \foreach \n in {0,1,2,3,4}
    {\draw[very thick, omDef] (0.5,{0.55+\n*1.05}) -- (3.3,{0.55+\n*1.05});
     \node[anchor=west] at (3.45,{0.55+\n*1.05})
       {$\ket{\n}$\; $E = (\n + \tfrac12)\hbar\omega$};}
  \draw[->, omThm, very thick] (1.35,0.62) -- (1.35,1.53);
  \node[omThm, anchor=east] at (1.28,1.08) {$\hat a^\dagger$};
  \draw[->, omExo, very thick] (2.45,1.53) -- (2.45,0.62);
  \node[omExo, anchor=west] at (2.52,1.08) {$\hat a$};
  \draw[<->] (0.15,0) -- (0.15,0.55);
  \node[anchor=west, font=\footnotesize] at (0.28,0.22) {$\tfrac12\hbar\omega$};
  \draw[dashed] (0,0) -- (3.3,0);
\end{tikzpicture}

{\small La escalera del oscilador: peldaños iguales $\hbar\omega$, subidos
por $\hat a^\dagger$ y bajados por $\hat a$, apoyada sobre un suelo situado
medio peldaño por encima de la energía de reposo clásica.}
\end{omfigure}

\begin{proposition}[Los estados en el espacio]\label{prop:b3:harmonic-oscillator:wavefunctions}
El estado fundamental es la gaussiana
\[
\varphi_0(x) = \Big(\frac{m\omega}{\pi\hbar}\Big)^{1/4}
\eu^{-m\omega x^2/2\hbar} ,
\]
obtenida al resolver la ecuación de \emph{primer orden} $\hat a\varphi_0
= 0$; satura la desigualdad de Heisenberg, $\Delta x\,\Delta p =
\hbar/2$, con $\Delta x = \sqrt{\hbar/2m\omega}$. Aplicar
$\hat a^\dagger$ repetidamente genera $\varphi_n$: un polinomio de
grado $n$ (con $n$ nodos) por la misma gaussiana. Para $n$ grande,
$|\varphi_n|^2$ oscila rápidamente en torno a la distribución clásica del
tiempo de permanencia, máxima cerca de los puntos de retorno --- el
principio de correspondencia en imagen.
\end{proposition}

\begin{proof}[Demostración parcial]
$\hat a\varphi_0 = 0$ se lee $\varphi_0' = -(m\omega/\hbar)x
\varphi_0$: la gaussiana, normalizada. La saturación: la gaussiana es
el caso de igualdad del \cref{thm:b3:quantum-formalism:uncertainty}.
La estructura polinómica se sigue de que $\hat a^\dagger$ es de primer
orden en $x$ y $\dd/\dd x$; el enunciado para $n$ grande se comprueba en el
\cref{exo:b3:harmonic-oscillator:12}.
\end{proof}

\begin{omfigure}
\begin{tikzpicture}
  \begin{axis}[omaxis, axis x line=bottom, axis y line=left,
    width=6.6cm, height=5.2cm, xmin=-4, xmax=4, ymin=0, ymax=3.6,
    xlabel={$x$ (en unidades de $\sqrt{\hbar/m\omega}$)}, ylabel={},
    xlabel style={at={(axis description cs:0.5,-0.1)}, anchor=north},
    xtick={-2,0,2}, ytick=\empty, clip=false]
    \addplot[omDef, thick, domain=-4:4, samples=90]
      {0.56*exp(-x*x)+0.25};
    \addplot[omThm, thick, domain=-4:4, samples=90]
      {1.13*x*x*exp(-x*x)+1.3};
    \addplot[omExo, thick, domain=-4:4, samples=120]
      {0.28*(2*x*x-1)^2*exp(-x*x)+2.35};
    \node[omDef, font=\scriptsize, anchor=west] at (axis cs:2.4,0.5) {$|\varphi_0|^2$};
    \node[omThm, font=\scriptsize, anchor=west] at (axis cs:2.4,1.55) {$|\varphi_1|^2$};
    \node[omExo, font=\scriptsize, anchor=west] at (axis cs:2.4,2.6) {$|\varphi_2|^2$};
  \end{axis}
  \begin{axis}[omaxis, axis x line=bottom, axis y line=left,
    width=6.6cm, height=5.2cm, xmin=-5.2, xmax=5.2, ymin=0, ymax=0.45,
    xshift=7.1cm,
    xlabel={$x$}, ylabel={},
    xlabel style={at={(axis description cs:0.5,-0.1)}, anchor=north},
    xtick=\empty, ytick=\empty, clip=false]
    \addplot[omDef, domain=-4.6:4.6, samples=300]
      {(64*x^6-480*x^4+720*x^2-120)^2*exp(-x*x)/81680};
    \addplot[omThm, very thick, dashed, domain=-3.54:3.54, samples=120]
      {0.3183/sqrt(max(0.35,13-x*x))};
    \node[omThm, font=\scriptsize, anchor=south] at (axis cs:0,0.36) {tiempo de permanencia clásico};
    \node[omDef, font=\scriptsize, anchor=west] at (axis cs:1.9,0.42) {$|\varphi_6|^2$};
  \end{axis}
\end{tikzpicture}

{\small Izquierda: las densidades de probabilidad más bajas (desplazadas
verticalmente): una gaussiana sin nodos y luego $n$ nodos para $\varphi_n$.
Derecha: para $n$ grande, la densidad cuántica oscila en torno a la
distribución clásica, que se acumula en los puntos de retorno, donde la masa
oscilante se demora.}
\end{omfigure}

\section{La energía del punto cero es real}

\begin{remark}[Por qué el suelo no puede estar más abajo]\label{rem:b3:harmonic-oscillator:zeropoint}
Un estado de energía nula exigiría $\langle\hat p^2\rangle =
\langle\hat x^2\rangle = 0$: perfectamente quieto \emph{y} perfectamente
centrado, algo que prohíbe $\Delta x\,\Delta p \ge \hbar/2$. El estado
fundamental es el mejor compromiso que la desigualdad permite, y
$\tfrac12\hbar\omega$ es el alquiler. No es una constante de contabilidad:
el movimiento de punto cero difumina los diagramas de difracción de rayos X
en el cero absoluto; da a los isótopos más ligeros enlaces efectivos más
débiles (H$_2$ y D$_2$ se disocian con energías medibles distintas a partir
del mismo pozo electrónico); y en el helio es tan violento --- átomos
ligeros, atracción endeble --- que el líquido \emph{nunca} se congela bajo
su propia presión de vapor: el único elemento todavía líquido en el cero
absoluto, que solo solidifica con la ayuda de 25 atmósferas.
\end{remark}

\begin{example}[Escalas de $\hbar\omega$]\label{ex:b3:harmonic-oscillator:scales}
Enlace del monóxido de carbono ($\omega = \qty{4.1e14}{rad/s}$): $\hbar\omega =
\qty{0.27}{eV}$, diez veces el $k_{\text{B}}T$ de la temperatura ambiente
--- las vibraciones moleculares están congeladas en el aire cotidiano, y por
eso las capacidades caloríficas de los gases diatómicos desconcertaron al
siglo XIX. Un péndulo ($\omega = \qty{5}{rad/s}$): $\hbar\omega = \qty{3e-15}{eV}$,
desesperadamente fuera de resolución --- el límite de correspondencia. Un
espejo de LIGO de \qty{40}{kg}, suspendido de modo que oscile a
\qty{100}{Hz}, tiene como modo de oscilador la amplitud de punto cero $\Delta x =
\sqrt{\hbar/2m\omega} \approx \qty{5e-20}{m}$ --- y el observatorio resuelve
de forma rutinaria desplazamientos \emph{en} este suelo cuántico: el
movimiento de punto cero de un objeto de cuarenta kilos es hoy una
restricción de ingeniería (\cref{exo:b3:harmonic-oscillator:3}).
\end{example}

\section{Estados coherentes: la cara clásica del oscilador cuántico}

\begin{definition}[Estados coherentes]\label{def:b3:harmonic-oscillator:coherent}
\index{estado coherente}
Un \emph{estado coherente} $\ket\alpha$ es un vector propio del \omterm{def:b3:harmonic-oscillator:ladder}{operador de
aniquilación}, $\hat a\ket\alpha = \alpha\ket\alpha$, con
$\alpha$ un número complejo cualquiera. Desarrollado sobre la escalera,
\[
\ket\alpha = \eu^{-|\alpha|^2/2}\sum_{n=0}^\infty
\frac{\alpha^n}{\sqrt{n!}}\,\ket n :
\]
el recuento cuántico sigue una distribución de Poisson de media $\langle\hat
N\rangle = |\alpha|^2$ y dispersión $\Delta N = |\alpha|$.
\end{definition}

\begin{proposition}[Por qué los láseres son clásicos]\label{prop:b3:harmonic-oscillator:coherent-motion}
Un \omterm{def:b3:harmonic-oscillator:coherent}{estado coherente} es la gaussiana del estado fundamental desplazada en el
\omterm{def:b3:hamiltonian-mechanics:hamiltonian}{espacio de fases}; bajo la evolución del oscilador \emph{sigue} siendo
coherente, con $\alpha(t) = \alpha\,\eu^{-\iu\omega t}$: su centro ejecuta
exactamente el movimiento clásico, $\langle\hat x\rangle(t) =
x_0\cos\omega t + \cdots$, mientras que sus anchuras permanecen en el mínimo
$\Delta x\,\Delta p = \hbar/2$ para siempre --- un paquete de ondas que
nunca se ensancha. Los \omterm{def:b3:harmonic-oscillator:coherent}{estados coherentes} son la manera en que un oscilador
cuántico se hace pasar por uno clásico; la luz de un láser es un \omterm{def:b3:harmonic-oscillator:coherent}{estado
coherente} de un modo del campo, con su número de fotones poissoniano (el
\emph{ruido de disparo} de todo fotodetector) y su campo oscilando como dijo
Maxwell.
\end{proposition}

\begin{proof}[Demostración parcial]
El desarrollo se sigue de escribir $\ket\alpha = \sum c_n\ket n$ en
$\hat a\ket\alpha = \alpha\ket\alpha$: $c_n = \alpha c_{n-1}/\sqrt
n$. Evolución: cada $\ket n$ adquiere $\eu^{-\iu(n + 1/2)\omega t}$,
que vuelve a sumarse en un \omterm{def:b3:harmonic-oscillator:coherent}{estado coherente} de $\alpha\eu^{-\iu\omega t}$
(salvo fase global). $\langle\hat x\rangle \propto
\operatorname{Re}\alpha(t)$ a partir de $\hat x \propto \hat a +
\hat a^\dagger$. Que el estado sea la gaussiana desplazada se admite aquí.
\end{proof}

\begin{omfigure}
\begin{tikzpicture}[font=\small, >=Stealth]
  \draw[->] (-3.1,0) -- (3.4,0) node[right] {$\langle\hat x\rangle$};
  \draw[->] (0,-3.0) -- (0,3.2) node[above] {$\langle\hat p\rangle$};
  \draw[gray, dashed] (0,0) circle (2.2);
  \fill[omDef!25, draw=omDef, thick] (2.2,0) circle (0.42);
  \fill[omDef!25, draw=omDef, thick] (50:2.2) circle (0.42);
  \fill[omDef!25, draw=omDef, thick] (100:2.2) circle (0.42);
  \draw[->, thick] (2.53,0.9) arc (20:40:2.6);
  \node[anchor=west] at (2.62,1.35) {$\omega$};
  \fill[omThm!30, draw=omThm, thick] (0,0) circle (0.42);
  \node[omThm, font=\footnotesize, anchor=east] at (-0.3,-0.68) {$\ket0$};
  \node[omDef, font=\footnotesize, anchor=west] at (2.3,-0.62) {$\ket\alpha$};
  \node[font=\footnotesize, anchor=west] at (-3.05,2.75)
    {mancha mínima $\Delta x\,\Delta p = \hbar/2$, girando clásicamente};
\end{tikzpicture}

{\small Retrato en el \omterm{def:b3:hamiltonian-mechanics:hamiltonian}{espacio de fases} (compárese con el
\cref{ch:b3:hamiltonian-mechanics}): el estado fundamental es una mancha de
incertidumbre mínima en el origen; un \omterm{def:b3:harmonic-oscillator:coherent}{estado coherente} es la misma mancha
desplazada, que gira a $\omega$ sin deformarse --- la mejor imitación de una
oscilación clásica que logra la mecánica cuántica.}
\end{omfigure}

\begin{method}[Álgebra del oscilador]\label{met:b3:harmonic-oscillator:method}
(1) Exprésese lo que se pida en $\hat a$ y $\hat a^\dagger$: $\hat x
= \sqrt{\hbar/2m\omega}\,(\hat a + \hat a^\dagger)$, $\hat p =
\iu\sqrt{m\hbar\omega/2}\,(\hat a^\dagger - \hat a)$. (2) Llévense los
$\hat a$ a la derecha con $[\hat a, \hat a^\dagger] = 1$; úsese
$\hat a\ket0 = 0$. (3) Elementos de matriz: $\hat x$ conecta solo peldaños
vecinos --- la regla de selección $\Delta n = \pm1$ de la espectroscopia
vibracional. (4) Todo \omterm{def:b3:hamiltonian-mechanics:hamiltonian}{hamiltoniano} cuadrático (osciladores acoplados,
circuitos, modos de campo) se diagonaliza en escaleras independientes:
hállense primero los \omterm{prop:b3:lagrangian-mechanics:modes}{modos normales} y cuantícese cada uno. (5) Números
primero: $\hbar\omega$ frente a $k_{\text{B}}T$ decide si el sistema es
cuántico o clásico antes de cualquier álgebra.
\end{method}

\section{Ejercicios}

\begin{exercise}[$\star$]\label{exo:b3:harmonic-oscillator:1}
(a) Verifíquese $[\hat a, \hat a^\dagger] = 1$ a partir de $[\hat x, \hat p] =
\iu\hbar$. (b) Verifíquese $\hat H = \hbar\omega(\hat a^\dagger\hat a +
\tfrac12)$ por desarrollo directo. (c) Inviértase para expresar $\hat x$ y
$\hat p$. (d) Muéstrese que $\hat N$ es hermítico y explíquese por qué
$\hat a$ por sí solo no podría ser un observable.
\end{exercise}

\begin{exercise}[$\star$]\label{exo:b3:harmonic-oscillator:2}
Sobre el estado propio $\ket n$: (a) muéstrese que $\langle\hat x\rangle =
\langle\hat p\rangle = 0$; (b) calcúlense $\langle\hat x^2\rangle$ y
$\langle\hat p^2\rangle$; (c) dedúzcase $\Delta x\,\Delta p = (n +
\tfrac12)\hbar$; (d) compruébese la relación del virial $\langle E_k\rangle =
\langle E_p\rangle$.
\end{exercise}

\begin{exercise}[$\star$]\label{exo:b3:harmonic-oscillator:3}
Amplitudes de punto cero $\Delta x = \sqrt{\hbar/2m\omega}$: calcúlense
para (a) la molécula de CO ($\mu = \qty{1.14e-26}{kg}$, $\omega =
\qty{4.1e14}{rad/s}$), comparada con la longitud de enlace de
\qty{0.11}{nm}; (b) un átomo de hidrógeno en un sólido ($m =
\qty{1.7e-27}{kg}$, $\hbar\omega = \qty{0.1}{eV}$); (c) un espejo de LIGO
($m = \qty{40}{kg}$, $f = \qty{100}{Hz}$); (d) ordénense las tres como
fracciones del tamaño de sus sistemas y coméntese quién es ``cuántico''.
\end{exercise}

\begin{exercise}[$\star$]\label{exo:b3:harmonic-oscillator:4}
(a) Usando las relaciones de escalera, calcúlese $\bra m\hat x\ket n$ y
muéstrese que se anula salvo si $m = n \pm 1$. (b) Dedúzcase la regla de
selección vibracional $\Delta n = \pm1$ para la absorción de luz (el
acoplamiento es $\propto\hat x$). (c) ¿Por qué absorbe luz infrarroja una
molécula heteronuclear (CO) y no la absorbe el N$_2$? (d) Las moléculas
reales muestran líneas débiles de ``sobretono'' en $\approx 2\hbar\omega$:
¿qué revela eso sobre el potencial?
\end{exercise}

\begin{exercise}[$\star\star$]\label{exo:b3:harmonic-oscillator:5}
(a) Resuélvase $\hat a\varphi_0 = 0$ como ecuación diferencial y
normalícese. (b) Genérense $\varphi_1$ y $\varphi_2$ aplicando
$\hat a^\dagger$. (c) Verifíquese $\varphi_1 \perp \varphi_0$ solo por
paridad. (d) Esbócense las tres densidades y compruébese el recuento de nodos.
\end{exercise}

\begin{exercise}[$\star\star$]\label{exo:b3:harmonic-oscillator:6}
Anticipo de Boltzmann. Un conjunto de osciladores idénticos a temperatura
$T$ ocupa el nivel $n$ con probabilidad $\propto
\eu^{-E_n/k_{\text{B}}T}$ (volumen del Año 2, factor de Boltzmann). (a)
Muéstrese que el número cuántico medio es $\langle n\rangle =
1/(\eu^{\hbar\omega/k_{\text{B}}T} - 1)$. (b) Evalúese para la vibración
del CO a \qty{300}{K} y a \qty{2000}{K}. (c) Muéstrese que la energía media
tiende a $k_{\text{B}}T$ a alta temperatura (equipartición recuperada) y a
$\tfrac12\hbar\omega$ a baja. (d) ¿A qué temperatura tiene una vibración de
\qty{15}{\micro m} (la flexión del dióxido de carbono) $\langle n\rangle = 0.1$?
\end{exercise}

\begin{exercise}[$\star\star$]\label{exo:b3:harmonic-oscillator:7}
Estadística de los \omterm{def:b3:harmonic-oscillator:coherent}{estados coherentes}. (a) A partir del desarrollo de
$\ket\alpha$, muéstrese que $\mathcal P(n)$ es de Poisson con media
$|\alpha|^2$. (b) Muéstrese que $\langle\hat N\rangle = |\alpha|^2$ y
$\Delta N = |\alpha|$. (c) Un láser de \qty{1}{mW} a \qty{633}{nm}: fotones
por segundo y fluctuación relativa $\Delta N/\langle N\rangle$ en un
segundo. (d) Ruido de disparo: muéstrese que la razón ruido/señal de la
fotocorriente cae como $1/\sqrt{\langle N\rangle}$ --- por qué los haces
intensos parecen suaves.
\end{exercise}

\begin{exercise}[$\star\star$]\label{exo:b3:harmonic-oscillator:8}
Evolución de un \omterm{def:b3:harmonic-oscillator:coherent}{estado coherente}. (a) Aplíquense las fases de evolución al
desarrollo y muéstrese que se obtiene $\ket{\alpha(t)}$ con $\alpha(t) =
\alpha\eu^{-\iu\omega t}$ (salvo una fase global). (b) Dedúzcanse
$\langle\hat x\rangle(t)$ y $\langle\hat p\rangle(t)$ y compárense
con la solución clásica. (c) ¿Por qué un estado propio de la energía, pese a
ser estacionario, \emph{no} describe un péndulo que oscila, y qué rasgo del
\omterm{def:b3:harmonic-oscillator:coherent}{estado coherente} sí lo hace? (d) Estímese $|\alpha|$ para un péndulo real
(\qty{10}{g}, \qty{10}{cm} de amplitud, \qty{1}{Hz}) y coméntese.
\end{exercise}

\begin{exercise}[$\star\star$]\label{exo:b3:harmonic-oscillator:9}
El circuito LC cuántico. Un lazo superconductor con $L =
\qty{10}{nH}$ y $C = \qty{0.4}{pF}$ hace oscilar la carga y el flujo
como $x$ y $p$. (a) Su frecuencia de resonancia (volumen del Año 1) y
el cuanto $\hbar\omega$ en $\unit{\micro eV}$. (b) ¿Por debajo de qué
temperatura se cumple $k_{\text{B}}T \ll \hbar\omega$, de modo que el
circuito quede en su estado fundamental? (c) ¿Por qué se operan los cúbits
superconductores en refrigeradores de dilución a \qty{20}{mK}? (d) La
fluctuación de tensión de punto cero $\Delta V = \Delta q/C$ con $\Delta q =
\sqrt{\hbar\omega C/2}$: evalúese.
\end{exercise}

\begin{exercise}[$\star\star\star$]\label{exo:b3:harmonic-oscillator:10}
Van der Waals a partir del movimiento de punto cero. Modelícense dos átomos
neutros a distancia $R$ como dos osciladores dipolares idénticos (carga $e$,
masa $m$, frecuencia $\omega_0$) cuyos desplazamientos se acoplan mediante
la energía dipolar $\lambda\,x_1x_2$ con $\lambda = e^2/2\pi\varepsilon_0R^3$.
(a) Muéstrese que las coordenadas normales $(x_1 \pm x_2)/\sqrt2$ oscilan a
$\omega_\pm = \omega_0\sqrt{1 \pm \lambda/m\omega_0^2}$. (b) La
energía fundamental es $\tfrac\hbar2(\omega_+ + \omega_-)$: desarróllese a
segundo orden en $\lambda$ y muéstrese que la energía de interacción es
\[
\Delta E = -\frac{\hbar\lambda^2}{8m^2\omega_0^3}
\ \propto\ -\frac{1}{R^6} .
\]
(c) ¿Por qué es universal esta atracción --- presente incluso entre átomos
sin dipolo permanente alguno? (d) La ley $1/R^6$ es la fuerza de van der
Waals de las correcciones al gas real del volumen del Año 1: ¿qué
``fluctúa'', microscópicamente, y qué le ocurriría a esta fuerza en un mundo
con $\hbar = 0$?
\end{exercise}

\begin{exercise}[$\star\star\star$]\label{exo:b3:harmonic-oscillator:11}
Bandas laterales de un ion atrapado. Un ion en una trampa armónica ($f =
\qty{1}{MHz}$) absorbe luz láser. Como el ion se mueve, su espectro de
absorción muestra la línea electrónica en $\nu_0$ flanqueada por
líneas en $\nu_0 \pm f$: transiciones que cambian en $\mp1$ el número
cuántico \emph{de movimiento} junto con el electrónico. (a) ¿Cuánto vale
$\hbar\omega_{\text{trampa}}$ en $\unit{neV}$ y por qué hace falta una línea
muy estrecha para resolver las bandas laterales? (b) Excitar la línea
$\nu_0 - f$ retira un cuanto de movimiento por ciclo: explíquese el
\emph{enfriamiento por banda lateral} hasta el estado fundamental. (c) Una
vez que $\langle n\rangle \approx
0$, la banda lateral inferior desaparece por completo --- ¿por qué es esa asimetría una prueba de haber alcanzado el
estado fundamental cuántico? (d) Así es como se certifica el estado
fundamental de movimiento de un solo átomo --- y, por otros medios, el de
los espejos de LIGO de escala kilogramo ---: dígase a qué ``temperatura'' ha
llegado el ion para $f = \qty{1}{MHz}$ y $\langle n\rangle = 0.05$.
\end{exercise}

\begin{exercise}[$\star\star\star$]\label{exo:b3:harmonic-oscillator:12}
El límite clásico, cuantitativamente. Un oscilador clásico de
amplitud $A$ pasa en $[x, x + \dd x]$ la fracción $\dd
t/T = \dd x/\pi\sqrt{A^2 - x^2}$. (a) Dedúzcase esta distribución del tiempo
de permanencia. (b) Para el estado cuántico $n$, tómese $A_n$ de $E_n =
\tfrac12 m\omega^2A_n^2$ y compárese la distribución clásica con
el $|\varphi_n|^2$ (dado) promediado localmente: ¿dónde coinciden y dónde
tienen que diferir? (c) Muéstrese que el espaciado relativo entre niveles
contiguos, $\Delta E/E$, se anula como $1/n$: la energía se vuelve
efectivamente continua. (d) Para el péndulo del
\cref{exo:b3:harmonic-oscillator:8}(d), estímense $n$ y $\Delta
E/E$, y conclúyase el argumento de correspondencia en una frase.
\end{exercise}

\section{Problema: la molécula que calienta la Tierra}

\begin{problem}[{Problema de fin de semana --- la escalera cuántica del
dióxido de carbono y el efecto invernadero}]\label{pb:b3:harmonic-oscillator:1}
Una molécula de dióxido de carbono es una cadena lineal O=C=O: unos pocos
osciladores cuantizados. Que el $\hbar\omega$ de su modo de flexión caiga
justo en mitad del resplandor térmico que la Tierra emite es la razón por la
que este gas traza --- cuatro moléculas de cada diez mil --- gobierna el
clima del planeta. Datos: número de ondas de la flexión $\tilde\nu_2 = \qty{667}{cm^{-1}}$
(la unidad del espectroscopista: $E = hc\tilde\nu$, $\qty{1}{cm^{-1}} =
\qty{1.24e-4}{eV}$); tensión asimétrica $\tilde\nu_3 =
\qty{2349}{cm^{-1}}$; tensión simétrica $\tilde\nu_1 =
\qty{1388}{cm^{-1}}$; $k_{\text{B}}T$ a $\qty{288}{K}$ vale
\qty{24.8}{meV}; ley de Wien (volumen del Año 2): $\lambda_{\max}T =
\qty{2898}{\micro m.K}$.

\medskip\noindent\textbf{Parte I --- Los modos de una molécula lineal.}
\begin{enumerate}
  \item Tres átomos tienen nueve \omterm{def:b3:lagrangian-mechanics:coordinates}{grados de libertad}: ¿cuántos son
        traslaciones de la molécula entera, cuántos rotaciones (la
        molécula es \emph{lineal}) y cuántas vibraciones quedan?
  \item Descríbanse las cuatro vibraciones: tensión simétrica, tensión
        asimétrica y una flexión doblemente degenerada --- ¿por qué
        aparece dos veces la flexión?
  \item La luz se acopla a un \emph{dipolo eléctrico} que vibra
        (\cref{exo:b3:harmonic-oscillator:4}). ¿Qué modos de
        O=C=O modulan el momento dipolar y cuál es mudo en el
        infrarrojo?
  \item Conviértanse los tres números de ondas en longitudes de onda de
        fotón y sitúense: ¿cuáles caen en el infrarrojo térmico?
  \item ¿Por qué no absorben prácticamente nada de infrarrojo los dos
        gases principales del aire, N$_2$ y O$_2$, y con qué consecuencia
        para la transparencia de la atmósfera?
  \item El vapor de agua, angular y dipolar, absorbe en buena parte del
        infrarrojo: ¿en qué ``ventana'' espectral opera casi en solitario
        la flexión del dióxido de carbono (compárense los \qty{15}{\micro m}
        con las bandas intensas del agua por debajo de \qty{8}{\micro m} y
        por encima de \qty{20}{\micro m})?
\end{enumerate}

\medskip\noindent\textbf{Parte II --- La escalera cuántica de la flexión.}
\begin{enumerate}[resume]
  \item Calcúlese $\hbar\omega_2$ en meV y la escalera $E_n$.
  \item ¿Qué fracción de moléculas ocupa $n = 1$ a
        \qty{288}{K} (respecto de $n = 0$, factor de Boltzmann)? ¿Y
        $n = 2$?
  \item ¿Qué longitud de onda de fotón impulsa $n = 0 \to 1$? ¿Por qué
        domina la \emph{misma} longitud de onda la emisión?
  \item Justifíquese, a partir de la escalera armónica, que una sola
        longitud de onda sirva a toda la escalera ($n \to n + 1$ para
        todo $n$): ¿qué propiedad del espaciado de niveles interviene?
  \item La molécula también rota, lo que añade estructura fina: el
        rasgo de \qty{15}{\micro m} es en realidad una \emph{banda} de
        aproximadamente \qty{1}{\micro m} de anchura. ¿Por qué importa la
        \emph{anchura} de banda para un gas de efecto invernadero (piénsese
        en qué ocurre una vez que el centro de la banda es opaco)?
  \item Un CO$_2$ vibracionalmente excitado en el aire tiene mucha más
        probabilidad de perder su cuanto por colisión que por radiación
        (vida media radiativa $\sim\qty{1}{s}$, tiempo entre colisiones
        $\sim\qty{e-9}{s}$): ¿adónde va realmente la energía radiante
        absorbida?
  \item A la inversa, el aire a \qty{288}{K} mantiene una población
        térmica en $n = 1$ (pregunta 8): ¿qué hace esa población que
        importe para el balance energético?
\end{enumerate}

\medskip\noindent\textbf{Parte III --- La contabilidad radiativa del planeta.}
\begin{enumerate}[resume]
  \item El Sol radia como un cuerpo a \qty{5800}{K}: calcúlese su pico de
        Wien. ¿Intercepta mucha luz solar la flexión del CO$_2$?
  \item El suelo radia como un cuerpo a \qty{288}{K}: calcúlese su pico
        de Wien y sitúense los \qty{15}{\micro m} en esa curva térmica.
  \item Explíquese el mecanismo de efecto invernadero en cuatro frases: la
        luz solar entra, el infrarrojo térmico sale, se intercepta en
        \qty{15}{\micro m} y se reemite tanto hacia arriba \emph{como}
        hacia abajo.
  \item La reemisión que escapa al espacio procede de capas altas y frías
        ($\sim\qty{220}{K}$): ¿por qué reduce la potencia saliente del
        planeta en esas longitudes de onda el hecho de emitir desde una
        capa más fría (recuérdese que la emisión térmica crece con $T$)?
  \item Más CO$_2$ empuja la capa emisora más arriba y más fría:
        enúnciese en una frase por qué debe entonces calentarse la
        superficie para reequilibrar las cuentas.
  \item El centro de la banda ya es opaco; el efecto del CO$_2$ añadido
        actúa en las \emph{alas} de la banda, lo que da un crecimiento
        logarítmico del forzamiento con la concentración: conéctese esto
        con la pregunta 11.
\end{enumerate}

\medskip\noindent\textbf{Parte IV --- Isótopos: la huella dactilar
cuántica.}
\begin{enumerate}[resume]
  \item La frecuencia de un oscilador va como $\sqrt{k/\mu}$: para la
        tensión asimétrica, sustituir $^{12}$C por $^{13}$C cambia la
        masa efectiva; la línea observada se desplaza de
        \qty{2349}{cm^{-1}} a unos \qty{2283}{cm^{-1}}. Compruébese el
        orden de magnitud de este desplazamiento del $\approx 3\%$ a
        partir de las masas.
  \item Unos láseres sintonizados con esas dos líneas cuentan por
        separado el $^{13}$CO$_2$ y el $^{12}$CO$_2$ de una muestra de
        gas: explíquese por qué la escalera cuantizada hace posible esa
        detección resuelta en isótopos.
  \item Las plantas prefieren el isótopo ligero, de modo que el carbono
        fósil es pobre en $^{13}$C: ¿qué ha hecho la razón isotópica
        medida del CO$_2$ atmosférico a medida que crecía su
        concentración, y qué demuestra eso sobre la \emph{fuente} del gas
        añadido?
  \item La misma física de los \qty{15}{\micro m} opera en Venus
        ($96\%$ de CO$_2$, \qty{90}{bar}): ¿qué ilustra su superficie a
        \qty{737}{K}?
  \item Marte respira también CO$_2$ casi puro, pero a \qty{6}{mbar},
        y es gélido: ¿qué demuestra el trío Venus--Tierra--Marte sobre
        qué variable controla la intensidad del efecto?
  \item Resúmase el resultado con nombre propio: un cuanto de
        \qty{83}{meV} --- el $\hbar\omega$ de la flexión de una molécula
        triatómica --- aparcado en \qty{15}{\micro m} sobre el espectro
        térmico de un planeta a \qty{288}{K}, absorbido, termalizado en
        nanosegundos y reemitido desde altitudes frías, es el mecanismo
        por el que un $0.04\%$ del aire fija la temperatura de la Tierra.
\end{enumerate}
\end{problem}
